\chapter{Methodology and Design}

\section{Overall Architecture}
The system is structured in layers:
\begin{enumerate}
    \item video acquisition,
    \item perception (detection + tracking + zone filtering),
    \item analysis (metrics + uncertainty + thresholds),
    \item output (visualization, CSV, webhook).
\end{enumerate}

\section{Processing Pipeline}
For each frame:
\begin{itemize}
    \item people detection,
    \item temporal association of trajectories,
    \item polygon-zone membership verification,
    \item update of entry/exit events,
    \item computation of queue indicators.
\end{itemize}

\section{Indicator Modeling}
\subsection{Arrival and Service Rates}
Observed events over a time window are used to estimate:
\begin{equation}
\lambda = \frac{N_{\text{arrivals}}}{\Delta t},
\qquad
\mu = \frac{N_{\text{exits}}}{\Delta t}.
\end{equation}

\subsection{Waiting Time}
Under the M/M/1 assumption, one form used in the implementation is:
\begin{equation}
W = \frac{\lambda}{\mu - \lambda}, \quad \text{with } \lambda < \mu.
\end{equation}

\section{Bayesian Uncertainty Quantification}
Rates are modeled using posterior Gamma distributions:
\begin{equation}
\lambda \sim \Gamma(\alpha_{\lambda}, \beta_{\lambda}),
\qquad
\mu \sim \Gamma(\alpha_{\mu}, \beta_{\mu}).
\end{equation}
Ninety-five percent credible intervals are computed for each metric.

\section{Threshold Detection and Alerts}
A dedicated module flags critical situations:
\begin{itemize}
    \item excessive waiting time,
    \item high queue backlog,
    \item arrival spikes,
    \item service degradation,
    \item instability or high uncertainty.
\end{itemize}

\section{Technology Choices}
\begin{itemize}
    \item Python, OpenCV, NumPy, Pandas,
    \item YOLO26 (Ultralytics) for detection,
    \item Supervision + ByteTrack for tracking,
    \item SciPy for statistical computation,
    \item Tkinter/CustomTkinter for the interface,
    \item Requests + n8n for external integration.
\end{itemize}

\section{Agile Project Management with Scrum}
\subsection{Product Backlog}
The product backlog was organized around business value and technical risk reduction. Priority user stories targeted: real-time queue visibility, reliable wait-time estimation, uncertainty-aware decision support, and external alerting through n8n/Telegram.

\subsection{Sprint Planning}
Planning was performed in short iterations, each with a clear sprint goal, implementation scope, and validation criteria. Work was decomposed into incremental deliverables to ensure traceable progress from detection pipeline to full integration.

\subsection{Scrum Ceremonies and Tracking}
Progress tracking relied on iterative review of completed features, pending tasks, and blockers. Sprint outcomes were documented in project logs, and backlog priorities were updated according to testing feedback and integration constraints.
